\chapter{Resta: primeros pasos}\label{ch:g1:subtraction}

Siete globos, dos se van volando: ¿cuántos quedan? Quitar es
\emph{restando}, escrito con el signo $-$. Y la resta tiene un
segunda cara: también responde "¿cuántos más?".

\section{quitando}

\begin{definition}[Sustracción]\label{def:g1:subtraction:def}
\emph{Restando} significa quitar y contar lo que queda. nosotros
escribir
\[
7 - 2 = 5
\]
y lea: ``siete menos dos es igual a cinco''.
\end{definition}

\begin{omfigure}
\begin{tikzpicture}[scale=0.6]
  \foreach \i in {0,1,2,3,4} \fill[omDef] (\i*1.05,0) circle (0.3);
  \foreach \i in {5,6} {
    \fill[omDef!30] (\i*1.05,0) circle (0.3);
    \draw[omThm, thick] (\i*1.05-0.32,-0.32) -- (\i*1.05+0.32,0.32);
    \draw[omThm, thick] (\i*1.05-0.32,0.32) -- (\i*1.05+0.32,-0.32);
  }
  \node[right, font=\small] at (7.6,0)
    {$7 - 2 = 5$: tacha dos, quedan cinco};
\end{tikzpicture}

{\small Una resta en imágenes: tachar es quitar.}
\end{omfigure}

\begin{example}[contando hacia atrás]\label{ex:g1:subtraction:countback}
Para una pequeña comida para llevar, cuente hacia atrás en la línea: $11 - 3$: decir
``once'', luego tres pasos atrás --- ``diez, nueve, ocho''. Entonces
$11 - 3 = 8$.
\end{example}

\section{¿Cuántos más?}

\begin{example}[la diferencia]\label{ex:g1:subtraction:difference}
``Nina tiene $9$ canicas y Tom tiene $6$. ¿Cuántos más tiene Nina?
No se quita nada, pero es una resta: empareja las canicas,
y contar las canicas de Nina sin pareja:
\[
9 - 6 = 3 .
\]
Nina tiene $3$ más. El resultado de una resta se llama
\emph{diferencia}\index{diferencia}.
\end{example}

\begin{method}[contando]\label{met:g1:subtraction:countup}
para encontrar un \omterm{ex:g1:subtraction:difference}{diferencia}, contando \emph{arriba} suele ser más fácil:
para $12 - 9$, empieza en $9$ y subir a $12$: ``diez, once,
doce'' --- tres pasos. Entonces $12 - 9 = 3$. (Subir a través $10$ para
brechas más grandes: $14 - 8$: de $8$ a $10$ es $2$, de $10$ a $14$
es $4$: \omterm{ex:g1:subtraction:difference}{diferencia} $6$.)
\end{method}

\begin{omfigure}
\begin{tikzpicture}[scale=0.6]
  \draw[-Stealth] (5.6,0) -- (15.4,0);
  \foreach \i in {6,7,8,9,10,11,12,13,14,15} \draw (\i,-0.1) -- (\i,0.1);
  \foreach \i in {8,10,14}
    \node[below, font=\footnotesize] at (\i,-0.12) {\i};
  \draw[-Stealth, omDef, thick] (8,0.3) .. controls (9,0.9) ..
    (10,0.3);
  \node[omDef, font=\small] at (9,1.0) {$+2$};
  \draw[-Stealth, omThm, thick] (10,0.3) .. controls (12,1.2) ..
    (14,0.3);
  \node[omThm, font=\small] at (12,1.25) {$+4$};
\end{tikzpicture}

{\small $14 - 8$ contando hacia arriba: desde $8$ a $14$ hay
$2 + 4 = 6$ pasos.}
\end{omfigure}

\section{La suma y la resta son gemelas.}

\begin{example}[familias de hechos]\label{ex:g1:subtraction:family}
los números $3$, $4$ y $7$ formar una pequeña familia:
\[
3 + 4 = 7, \qquad
4 + 3 = 7, \qquad
7 - 4 = 3, \qquad
7 - 3 = 4 .
\]
Conocer un dato de la familia le da los otros tres de forma gratuita.
La resta deshace la \omterm{def:g1:addition:def}{suma}: después $+4$, haciendo $-4$ te trae de vuelta.
\end{example}

\begin{example}[De cheques]\label{ex:g1:subtraction:check}
Calculado $13 - 5 = 8$? Verifique con la \omterm{def:g1:addition:def}{suma} de gemelos:
$8 + 5 = 13$. \checkmark\ Una resta es correcta exactamente al sumar
La espalda funciona.
\end{example}

\begin{example}[¿Qué operación?]\label{ex:g1:subtraction:choose}
``$6$ pájaros en el alambre, $4$ volar'': llevar, $6 - 4 = 2$.
``$6$ pájaros, entonces $4$ llegan más '': juntos, $6 + 4 = 10$.
``$6$ gorriones y $4$ petirrojos: ¿cuántos gorriones más?'':
\omterm{ex:g1:subtraction:difference}{diferencia}, $6 - 4 = 2$. ¡Dibuja la historia cuando no estés seguro!
\end{example}

\section{Ejercicios}

\begin{exercise}[$\star $]\label{exo:g1:subtraction:1}
Dibuja, tacha y calcula: $6 - 2$; \quad $8 - 5$; \quad
$10 - 4$.
\end{exercise}

\begin{exercise}[$\star $]\label{exo:g1:subtraction:2}
Cuente hacia atrás para calcular: $9 - 3$; \quad $12 - 2$; \quad
$15 - 4$.
\end{exercise}

\begin{exercise}[$\star $]\label{exo:g1:subtraction:3}
Cuenta hacia arriba para encontrar el \omterm{ex:g1:subtraction:difference}{diferencia}: $11 - 8$; \quad $13 - 9$; \quad
$16 - 12$.
\end{exercise}

\begin{exercise}[$\star $]\label{exo:g1:subtraction:4}
Escribe los cuatro datos de la familia de $2$, $6$ y $8$.
\end{exercise}

\begin{exercise}[$\star $]\label{exo:g1:subtraction:5}
Calcule y luego verifique con la \omterm{def:g1:addition:def}{suma} de gemelos: $14 - 6$; \quad
$17 - 9$.
\end{exercise}

\begin{exercise}[$\star $]\label{exo:g1:subtraction:6}
Completa los agujeros:
\[
10 - \;?\; = 7, \qquad
\;?\; - 5 = 5, \qquad
12 - \;?\; = 12 .
\]
\end{exercise}

\begin{exercise}[$\star $]\label{exo:g1:subtraction:7}
``Hay $15$ cerezas en el bol. los niños comen $8$. como
¿Cuántos quedan?'' Escribe la resta y la oración de respuesta.
\end{exercise}

\begin{exercise}[$\star $]\label{exo:g1:subtraction:8}
Para cada historia, elija $+$ o $-$, luego calcula: ``$9$ patos en el
estanque, $5$ alejarse nadando''; ``$9$ patos, $5$ llegan los gansos'';
``Alí es $9$, su hermana es $5$: ¿Cuántos años mayor tiene Ali?''.
\end{exercise}

\begin{exercise}[$\star\star $]\label{exo:g1:subtraction:9}
madre pato tenia $12$ huevos. Algunos nacieron; ahora $12 - ?$ quedan huevos
y ella cuenta $7$ huevos sin eclosionar. ¿Cuántos huevos eclosionaron? (¿Cuál
¿De hecho la familia es esta?)
\end{exercise}
